\thispagestyle{plain}
\begin{center}
    {\fontsize{14pt}{18pt}\selectfont \textbf{EXECUTIVE SUMMARY}\par}
\end{center}
\vspace{0.8cm}

\noindent Rapid, unplanned urbanization and recurrent high-intensity monsoon precipitation severely overwhelm aging stormwater drainage infrastructures across tier-two Indian municipalities, culminating in systemic waterlogging, structural asset degradation, and public paralysis. Conventional municipal public works operate reactively, deploying desilting crews primarily in response to localized citizen grievances or visible surface overflow while neglecting topological connectivity and backwater surcharge dynamics. 

This dissertation presents \textbf{CivicTwin}, an AI-enabled cyber-physical digital twin developed for predictive urban drainage management and proactive flood resilience in Vellore Municipal Corporation (CMC). CivicTwin formulates 13,998 drain conduits across 60 spatial zones into a Directed Acyclic Graph (DAG) comprising 22,737 hydraulic edges. Integrating the rational runoff method, an ensemble Random Forest blockage surrogate, and a bitset-accelerated reachability engine, CivicTwin introduces a multi-criteria cascade-aware dispatch strategy. Across simulated rainfall intensities from 80 to 160~mm/hr, the proposed cascade-aware method resolves between $3.1\times$ and $4.4\times$ more at-risk and backwater-threatened conduits than traditional severity-only desilting under identical municipal clearing budgets. Implemented as an interactive web GIS command dashboard connecting FastAPI, React-Leaflet, and PostgreSQL/PostGIS, CivicTwin provides municipal administrators with real-time hydraulic bottleneck foresight, maximizing labor productivity and transforming municipal disaster response into predictive urban resilience.

\clearpage
